% =====================================================================
% SEC_00 — Capa, folha de rosto, resumo, sumário e listas
% =====================================================================

% Capa
\imprimircapa

% Folha de rosto
\newpage
\begin{center}
{\large\imprimirtitulo}
\end{center}
\vspace{1cm}
\begin{center}
{\large Autor: \imprimirautor}
\end{center}
\vspace{0.2cm}
\begin{center}
{\small ORCID: \texttt{0000-0001-8996-2887} \quad E-mail: \texttt{marceloclaro@gmail.com}}
\end{center}
\vspace{0.5cm}
\begin{center}
{\normalsize Especificação: SPEC-935-R456 (Rodada R456)}
\end{center}
\vspace{1cm}

\textbf{Resumo:} Este manual documenta a arquitetura de Recuperação Aumentada por
Geração (RAG) do ecossistema, distinguindo rigorosamente o \textbf{estado atual}
(implementado e observável no código-fonte) da \textbf{arquitetura-alvo proposta}
(não implementada) que integra um Diversificador Estruturado baseado na sequência
de Recamán para aumentar a diversidade semântico-estrutural dos resultados
recuperados. O documento parte de um \textbf{problema de pesquisa} formalizado,
que motiva o estudo de diversificação estruturada, e apresenta mapas de
arquitetura, diagramas com legendas, cálculos e especificações técnicas
reproduzíveis, apoiados em referências reais verificadas (DOIs e preprints arXiv).
Foi escrito para uma trilha de leitura de profundidade progressiva --- do leitor
de nível introdutório ao leitor de pós-graduação.

\vspace{0.5cm}
\textbf{Palavras-chave:} RAG; Recuperação de Informação; Recamán; Diversificação;
Geração Aumentada por Recuperação; Arquitetura de Software; Problema de Pesquisa.

\newpage

% ---------------------------------------------------------------
% Sumário
% ---------------------------------------------------------------
\pdfbookmark[0]{\contentsname}{toc}
\tableofcontents*
\clearpage

% Lista de figuras e tabelas (formato banca)
\pdfbookmark[0]{\listfigurename}{lof}
\listoffigures*
\clearpage
\pdfbookmark[0]{\listtablename}{lot}
\listoftables*
\clearpage
